\chapter{Теория}

\section{Используемые инструменты}
В ходе работы использовалась такая библиотека как Scikit-learn. \\

\subsection{Scikit-learn (sklearn)}
Scikit-learn - Python-библиотека с открытым исходным кодом для классического машинного 
обучения, обеспечивающая эффективные инструменты для анализа данных, классификации, 
регрессии, кластеризации и снижения размерности.

\section{Метод ближайших соседей Knn}

К-ближайших соседей (K-Nearest Neighbors или просто KNN) — алгоритм классификации и регрессии, основанный на гипотезе компактности,
которая предполагает, что расположенные близко друг к другу объекты в пространстве признаков имеют схожие значения целевой 
переменной или принадлежат к одному классу.


Алгоритм строится следующим образом:

\begin{enumerate}
    \item Сначала вычисляется расстояние между тестовым и всеми обучающими образцами;
    \item Далее из них выбирается k-ближайших образцов (соседей), где число k задаётся заранее;
    \item Итоговым прогнозом среди выбранных k-ближайших образцов будет мода в случае классификации и среднее арифметическое в случае регрессии;
    \item Предыдущие шаги повторяются для всех тестовых образцов.
\end{enumerate}



\endinput